\documentclass[tikz,border=4pt]{standalone}
\usepackage{pe-tikz}
\begin{document}
\begin{tikzpicture}[x=1.2cm, y=0.62cm, tr/.style={fill=cA!35, draw=white, line width=1pt}, te/.style={fill=cB, draw=white, line width=1pt}]
  \node[font=\small\bfseries, text=peblue, anchor=west] at (0,6.2) {(a) 5-szörös keresztvalidáció};
  \foreach \r in {1,...,5}{
    \foreach \c in {1,...,5}{
      \ifnum\c=\r \fill[te] (\c-1,6-\r-0.4) rectangle (\c,6-\r+0.4); \else \fill[tr] (\c-1,6-\r-0.4) rectangle (\c,6-\r+0.4); \fi}
    \node[pe/note, anchor=east] at (-0.1,6-\r) {\r. kör};
    \node[pe/note, anchor=west] at (5.1,6-\r) {$E_\r$};}
  \node[pe/label, anchor=west] at (5.6,3) {$\hat{E}_{\mathrm{CV}} = \frac{1}{5}\sum_{j=1}^{5}E_j$};
  \begin{scope}[xshift=10.4cm]
  \node[font=\small\bfseries, text=peblue, anchor=west] at (0,6.2) {(b) idősoros (gördülő kezdetű) validáció};
  \foreach \r/\n in {1/2, 2/3, 3/4, 4/5}{
    \foreach \c in {1,...,\n}{\fill[tr] (\c-1,6-\r-0.4) rectangle (\c,6-\r+0.4);}
    \pgfmathtruncatemacro{\t}{\n+1}
    \fill[te] (\n,6-\r-0.4) rectangle (\t,6-\r+0.4);
    \node[pe/note, anchor=east] at (-0.1,6-\r) {\r. kör};}
  \draw[pe/thin, draw=muted] (0,0.6) -- (6.3,0.6) node[right, pe/note]{idő};
  \end{scope}
  \fill[tr] (0,-0.6) rectangle (0.6,-0.2); \node[pe/note, anchor=west] at (0.7,-0.4) {tanítás};
  \fill[te] (2.4,-0.6) rectangle (3.0,-0.2); \node[pe/note, anchor=west] at (3.1,-0.4) {validálás};
\end{tikzpicture}
\end{document}
